\section*{Chapter \ref{ch:g1:landscapes-around-us} --- Landscapes Around Us}

\begin{solution}{exo:g1:landscapes-around-us:1}
\omterm{def:g1:landscapes-around-us:landform}{Mountain}, \omterm{def:g1:landscapes-around-us:landform}{hill}, \omterm{def:g1:landscapes-around-us:landform}{valley}, \omterm{def:g1:landscapes-around-us:landform}{plain}.
\end{solution}

\begin{solution}{exo:g1:landscapes-around-us:2}
A \omterm{def:g1:landscapes-around-us:landform}{mountain}: it is higher and steeper than a \omterm{def:g1:landscapes-around-us:landform}{hill}.
\end{solution}

\begin{solution}{exo:g1:landscapes-around-us:3}
Natural: the \omterm{def:g1:landscapes-around-us:water}{river}, the forest, the beach. Built: the road, the bridge, the
house.
\end{solution}

\begin{solution}{exo:g1:landscapes-around-us:4}
Where the land meets the sea.
\end{solution}

\begin{solution}{exo:g1:landscapes-around-us:5}
For example: natural, the \omterm{def:g1:landscapes-around-us:landform}{mountains}, the forest on the slopes or the
stream; built, a house or the church of the village.
\end{solution}

\begin{solution}{exo:g1:landscapes-around-us:6}
The \omterm{def:g1:landscapes-around-us:landform}{mountain}, with snow on its top. The \omterm{def:g1:landscapes-around-us:water}{river} flows at the bottom of the
\omterm{def:g1:landscapes-around-us:landform}{valley}, between the \omterm{def:g1:landscapes-around-us:landform}{mountains} and the \omterm{def:g1:landscapes-around-us:landform}{hill}.
\end{solution}

\begin{solution}{exo:g1:landscapes-around-us:7}
A \omterm{def:g1:landscapes-around-us:water}{river} flows from higher land (\omterm{def:g1:landscapes-around-us:landform}{hills} and \omterm{def:g1:landscapes-around-us:landform}{mountains}) down to a \omterm{def:g1:landscapes-around-us:water}{lake} or to
the sea. Water always runs downhill, so it cannot flow up to the top of a
\omterm{def:g1:landscapes-around-us:landform}{hill}.
\end{solution}

\begin{solution}{exo:g1:landscapes-around-us:8}
The tree is near us, in the foreground; the \omterm{def:g1:landscapes-around-us:landform}{hill} is far away, in the
background. Things that are far away look smaller than things that are
near.
\end{solution}

\begin{solution}{exo:g1:landscapes-around-us:9}
Foreground: tall grass and a few bushes at the bottom of the picture.
Middle ground: the \omterm{def:g1:landscapes-around-us:water}{river} winding across the \omterm{def:g1:landscapes-around-us:landform}{plain} with trees along it,
fields, a farm and a road.
Background: the flat land going on to the horizon, and the sky.
\end{solution}

\begin{solution}{exo:g1:landscapes-around-us:10}
Because people made it: they cleared the land, ploughed it and sowed the
wheat. Without people there would be grass, bushes or forest there, not a
field of wheat.
\end{solution}

\begin{solution}{exo:g1:landscapes-around-us:11}
Not changed: the \omterm{def:g1:landscapes-around-us:landform}{hill} and the \omterm{def:g1:landscapes-around-us:water}{river}. Added by people: the fields, more
houses, the road and the bridge (and most of the forest has been cut).
\end{solution}
